\subsection{Known Limitations}

\begin{table}[H]
  \centering
  \caption{Known limitations of release 1.0.0}
  \label{tab:limitations}
  \begin{tabularx}{\textwidth}{l L}
    \toprule
    \textbf{Area} & \textbf{Limitation} \\
    \midrule
    Scope & Momentary current only. There is no duty-cycle or ampere-hour battery sizing. \\
    Protection logic & Zones are proposed from the drawing's topology. The engineer enters and verifies relay assignments, trip logic and operating sequences; they are not read from relay settings. \\
    Voltage & Currents are not adjusted for DC voltage. Reports state a nominal \SI{125}{V} system. \\
    PDF background & Only the first page of a PDF is used. The background is an image, so symbols are not recognized from it. \\
    Maintenance status & For reference only; a device in maintenance is treated as closed for zone detection. \\
    Portable package & Windows x64 only. It is not code-signed, so organizations may need to approve it. \\
    Testing & Mouse interaction on the drawing is tested by hand, not by automated tests. \\
    \bottomrule
  \end{tabularx}
\end{table}

\subsection{Future Improvements}

The most valuable addition would be \textbf{automated diagram recognition}: using
artificial intelligence (AI) to identify the components, symbols, labels and
connections in an imported PDF and turn them into editable equipment. Rebuilding the
drawing by hand is currently the slowest step of a study. As with detected zones
today, the engineer would still review the recognized equipment.

Other possible improvements:

\begin{itemize}
  \item \textbf{Battery sizing.} Extend the momentary results into a duty-cycle
        analysis following IEEE Std~485~\cite{ieee485}, using each supply's continuous,
        momentary and random loads.
  \item \textbf{Relay and scheme import.} Read trip assignments from relay settings or
        scheme drawings instead of entering them, reducing manual work and errors.
  \item \textbf{Manufacturer data library.} A shared library of breakers and motor
        operators with their published coil and motor currents, so that each current is
        entered only once.
  \item \textbf{Richer PDF support.} Choose a page from a multi-page drawing set, and
        align or scale the background interactively.
  \item \textbf{Equipment library.} Additional symbols, such as protective relays,
        capacitor banks and reactors, with their own properties.
  \item \textbf{Report options.} Export results to a spreadsheet, and compare two
        versions of a study to show what changed.
  \item \textbf{Deployment.} Sign the Windows package so organizations can approve it
        more easily, and automate the testing of mouse interaction on the drawing.
\end{itemize}
